% Архив прежнего текста; действующая методичка собирается из demos/seminar04/README.md.
\label{seminar04}
\begingroup
\setlength{\emergencystretch}{2em}
\subsection*{Цель, время и порядок работы}
За 90 минут студент на своём актуальном Debian установит инструменты,
самостоятельно соберёт и запустит односекторную BIOS-программу, затем
соберёт готовый журнал на 16-битном NASM и исследует вывод, ввод,
вычисления и сохранение в отдельный диск. Дома на это выделены ещё
180 минут. Студент выполняет команды терминала сам; проверяет экран,
вывод и образ диска своими наблюдениями и записывает их в отчёте.
Команды ниже выполняются из рабочего каталога, который создаётся
внутри репозитория. Для лабораторной не нужны Docker, make или
автоматические тесты; требуются Debian с графической средой и доступом
к \texttt{sudo}. Проверки материалов преподавателем в контейнерах
не являются частью действий студента.

\subsection*{В аудитории: 90 минут}
\paragraph{0--10 минут: пакеты Debian.}
В терминале установите необходимые программы:
\begin{lstlisting}[basicstyle=\ttfamily\small]
sudo apt update
sudo apt install qemu-system-x86 qemu-system-gui \
    nasm gcc gdb binutils
nasm -v
qemu-system-i386 --version
gcc --version
gdb --version
objdump --version
\end{lstlisting}
\texttt{apt update} обновляет сведения о пакетах, \texttt{apt install}
устанавливает ассемблер NASM, эмулятор ПК QEMU с окном GTK,
компилятор C, отладчик GDB и binutils. Флаги \texttt{--version} и
\texttt{-v} проверяют, что команды доступны; точная версия зависит
от Debian. Для просмотра машинных байтов понадобятся также штатные
\texttt{wc}, \texttt{od}, \texttt{dd}, \texttt{cp}, \texttt{truncate}.

\paragraph{10--25 минут: первый загрузочный сектор.}
Создайте рабочий каталог, не затирая существующую копию исходника:
\begin{lstlisting}[basicstyle=\ttfamily\small]
mkdir -p work/seminar04
cp -n demos/seminar04/hello.asm work/seminar04/hello.asm
cd work/seminar04
\end{lstlisting}
Прочитайте \texttt{hello.asm}: \texttt{bits 16}, \texttt{org 0x7c00},
установку \texttt{DS} и \texttt{SS:SP}, цикл \texttt{lodsb},
BIOS \texttt{INT 10h}, дополнение \texttt{times} и подпись сектора.
Соберите плоский бинарный файл, посмотрите его длину и последние байты:
\begin{lstlisting}[basicstyle=\ttfamily\small]
nasm -f bin hello.asm -o hello.bin
wc -c hello.bin
od -An -tx1 -j 510 -N 2 hello.bin
ndisasm -b 16 hello.bin
cp hello.bin hello.img
truncate -s 1440K hello.img
wc -c hello.img
\end{lstlisting}
В NASM \texttt{-f bin} означает плоские машинные байты без заголовка ELF,
\texttt{-o} задаёт выходной файл. \texttt{wc -c} считает байты:
загрузочный сектор должен иметь размер 512 байтов.
У \texttt{od} параметр \texttt{-An} убирает колонку адресов,
\texttt{-tx1} показывает каждый байт в hex, \texttt{-j 510}
пропускает 510 байтов, \texttt{-N 2} читает два байта подписи
\texttt{55 AA}. Команда \texttt{ndisasm -b 16} читает байты как
16-битные инструкции, но не знает, где закончился выполняемый код:
строку, нули и подпись она тоже может показать как команды.
\texttt{truncate -s 1440K} доводит образ дискеты
до $1440\cdot1024$ байтов.

Запустите этот файл в QEMU:
\begin{lstlisting}[basicstyle=\ttfamily\footnotesize]
qemu-system-i386 -machine pc -m 16M -boot order=a \
  -drive file=hello.img,format=raw,if=floppy \
  -display gtk -no-reboot
\end{lstlisting}
\texttt{-machine pc} выбирает BIOS-совместимый ПК, \texttt{-m 16M}
выделяет 16 MiB RAM гостю, \texttt{-boot order=a} выбирает дискету
первой для загрузки. У \texttt{-drive} параметры \texttt{file},
\texttt{format=raw}, \texttt{if=floppy} задают файл без заголовка
контейнера и его интерфейс. \texttt{-display gtk} открывает отдельное
окно, \texttt{-no-reboot} запрещает автоматический перезапуск гостя.
Проверьте надпись \texttt{HELLO BIOS}, поменяйте сообщение на короткое
латинское, повторите сборку и запуск и запишите наблюдения.
Этот бинарник не запускают командой \texttt{./hello.bin} в Linux.

\paragraph{25--35 минут: полный журнал.}
Из рабочего каталога вернитесь в корень репозитория за исходниками,
затем снова войдите в рабочий каталог:
\begin{lstlisting}[basicstyle=\ttfamily\footnotesize]
cd ../..
cp -n demos/lecture04-journal/boot16.asm work/seminar04/boot16.asm
cp -n demos/lecture04-journal/stages/stage6/journal.asm \
  work/seminar04/journal.asm
cd work/seminar04
nasm -f bin boot16.asm -o boot.bin
nasm -f bin journal.asm -o journal.bin
wc -c boot.bin journal.bin
\end{lstlisting}
\texttt{boot.bin} должен занимать 512 байтов; код журнала должен
уместиться в следующие 17 секторов по 512 байтов каждый. Найдите
максимальную длину самостоятельно. Сборка дискеты тремя действиями:
\begin{lstlisting}[basicstyle=\ttfamily\footnotesize]
dd if=/dev/zero of=boot.img bs=512 count=2880
dd if=boot.bin of=boot.img bs=512 conv=notrunc
dd if=journal.bin of=boot.img bs=512 seek=1 conv=notrunc
wc -c boot.img
\end{lstlisting}
\texttt{if=} и \texttt{of=} --- источник и результат,
\texttt{bs=512} задаёт размер блока, \texttt{count=2880}
задаёт число блоков на дискете объёмом 1,44 MiB. У третьего
\texttt{dd} параметр \texttt{seek=1} пропускает \emph{один блок
выходного файла}: журнал начинается со второго сектора.
\texttt{conv=notrunc} сохраняет оставшуюся часть образа.

Создайте \emph{один раз} отдельный HDD и запустите журнал:
\begin{lstlisting}[basicstyle=\ttfamily\footnotesize]
truncate -s 1M grades.img
qemu-system-i386 -machine pc -m 16M -boot order=a \
  -drive file=boot.img,format=raw,if=floppy \
  -drive file=grades.img,format=raw,if=ide,index=0,media=disk \
  -display gtk -serial stdio -monitor none -no-reboot
\end{lstlisting}
Второй \texttt{-drive} подключает \emph{отдельный} диск:
\texttt{if=ide} задаёт интерфейс IDE, \texttt{index=0} выбирает
первый накопитель, \texttt{media=disk} отличает HDD от дискеты.
Для BIOS это диск \texttt{80h}. \texttt{-serial stdio} выводит COM1
в терминал, \texttt{-monitor none} отключает консоль команд QEMU.
Проверьте исходное поведение: 0, стрелка вправо, A, S,
\texttt{SAVED}; завершите QEMU, запустите той же командой,
увидьте \texttt{LOADED}. После изменения программы повторите NASM
для \texttt{journal.bin} и все три команды \texttt{dd} для чистой
дискеты: старые байты после сокращённой программы не должны оставаться
в образе. Не создавайте заново \texttt{grades.img}, иначе сравнение
сохранений потеряет смысл.

\paragraph{35--50 минут: вывод, цвет и позиция.}
В \texttt{journal.asm}, процедуре \texttt{vga\_text}, найдите
\texttt{mov ah,0x0f} и сделайте текст жёлтым на синем фоне без мигания.
Таблица кодов текстового VGA:
\begin{center}\small
\begin{tabular}{|c|l|c|l|}\hline
Код & Цвет & Код & Цвет \\ \hline
0 & чёрный & 8 & тёмно-серый \\
1 & синий & 9 & ярко-синий \\
2 & зелёный & A & ярко-зелёный \\
3 & бирюзовый & B & ярко-бирюзовый \\
4 & красный & C & ярко-красный \\
5 & пурпурный & D & ярко-пурпурный \\
6 & коричневый & E & жёлтый \\
7 & светло-серый & F & белый \\ \hline
\end{tabular}\end{center}
В атрибуте бит 7 --- мигание (здесь 0), биты 6--4 --- фон
(0--7), биты 3--0 --- цвет символа (0--F). Вычислите
$16\cdot\text{код фона}+\text{код символа}$ сами. Одна позиция
VGA состоит из байта символа и байта атрибута.

В \texttt{redraw} найдите \texttt{mov bx,0x0001} перед
\texttt{mov si,title\_text}. Перенесите заголовок на строку 2,
столбец 5. Координаты начинаются с нуля, \texttt{BH} обозначает
строку, \texttt{BL} --- столбец. Рассчитайте новую константу
\texttt{BX} $=256\cdot\text{строка}+\text{столбец}$, смещение
$2\cdot(80\cdot\text{строка}+\text{столбец})$ и физический адрес
первой буквы от \texttt{B800:0000}. После сборки проверьте положение
и цвет в окне QEMU.

\paragraph{50--60 минут: ввод клавиш.}
В \texttt{handle\_scan}, ветке \texttt{.normal}, найдите сравнения
перед \texttt{.save} и \texttt{.load}. Назначьте Q для сохранения,
E для загрузки вместо S/L; обновите строку \texttt{help\_text}.
Сравниваются скан-коды PS/2 set 1, \emph{не ASCII}. Их десятичные
значения: S --- 31, L --- 38, Q --- 16, E --- 18, A --- 30.
Переведите новые коды в hex самостоятельно. Проверьте новые и старые
клавиши по строкам \texttt{SAVED}/\texttt{LOADED} в терминале.

\paragraph{60--75 минут: вычисления.}
Клавиша A должна выставлять 12 вместо 10. Найдите в
\texttt{handle\_scan} метку \texttt{.ten}, а в
\texttt{load\_grades} метку \texttt{.check}; поменяйте выдаваемую
оценку и максимальную допустимую оценку при загрузке. Обновите
подсказку. Делитель 10 в \texttt{vga\_two} сохраняется: он делит
число на десятичные разряды, а не проверяет предел оценки.
Рассчитайте вручную средние $12,0$ и $12,1,1$ с округлением до
сотых, проверьте на экране и после сохранения и перезапуска.

\paragraph{75--85 минут: новый сектор HDD.}
В \texttt{dap} найдите поле \texttt{dq 1}, номер LBA. Требуется
запись с байта 4096 отдельного HDD; сектор имеет 512 байтов.
Вычислите $LBA=4096/512$ и измените поле. Сохраните старый диск
как копию и запустите опыт с новым файлом:
\begin{lstlisting}[basicstyle=\ttfamily\footnotesize]
cp -n grades.img grades-lab.img
qemu-system-i386 -machine pc -m 16M -boot order=a \
  -drive file=boot.img,format=raw,if=floppy \
  -drive file=grades-lab.img,format=raw,if=ide,index=0,media=disk \
  -display gtk -serial stdio -monitor none -no-reboot
\end{lstlisting}
В новом секторе изначально может быть \texttt{EMPTY}. Выставьте
12 и 0, сохраните Q, повторите запуск с \emph{той же копией} и
увидьте \texttt{LOADED}. Найдите \texttt{GR16} с помощью
\texttt{od -An -tc -j} и самостоятельно найденного смещения;
\texttt{-tc} показывает четыре прочитанных байта как символы.

\paragraph{85--90 минут: устный разбор.}
Объясните цепочку «клавиша $\to$ числовая оценка в RAM $\to$
символы и атрибут VGA $\to$ секторный буфер RAM $\to$ HDD».
Чем отличаются \texttt{1000:0000}, \texttt{B800:...} и
\texttt{1000:sector}? Результат каждого действия проверяйте
своими наблюдениями, а не ответом автоматического теста.

\subsection*{Самостоятельная работа: 180 минут}
\paragraph{35 минут: размеры и байты в C/GDB.}
В рабочем каталоге выполните:
\begin{lstlisting}[basicstyle=\ttfamily\footnotesize]
gcc -std=c11 -Wall -Wextra -Werror -g -O0 \
  ../../demos/seminar04/memory.c -o memory
./memory
gdb -q ./memory
\end{lstlisting}
Флаг \texttt{-std=c11} задаёт стандарт языка;
\texttt{-Wall -Wextra -Werror} включает диагностику;
\texttt{-g} добавляет сведения для GDB, \texttt{-O0} отключает
оптимизации, \texttt{-o memory} задаёт имя файла,
\texttt{-q} выключает заставку GDB. \emph{Вручную} введите
\texttt{break show\_memory}, \texttt{run},
\texttt{print sizeof(char)}, \texttt{print sizeof(int)},
\texttt{print sizeof(long)}, \texttt{print *x},
\texttt{x/4xb x}, \texttt{x/3dw numbers}, затем
\texttt{continue} и \texttt{quit}. В команде \texttt{x/4xb}
число 4 задаёт количество, \texttt{x} --- hex, \texttt{b} --- байт;
\texttt{x/3dw} читает три 32-битных слова в десятичном виде.
Сравните размеры типов, шаги между элементами массивов и байты
\texttt{0x12345678}. Виртуальные адреса Linux и физическая память
журнала --- разные адресные пространства.

\paragraph{35 минут: цвета и адреса.}
Вычислите атрибуты «ярко-зелёный на чёрном» и «белый на красном»,
потом \texttt{BX}, смещение и физические адреса для ещё двух
свободных позиций текстового экрана. Укажите системы счисления.

\paragraph{30 минут: ввод и арифметика.}
Переведите скан-коды Q/E в hex; рассчитайте средние $12,0$ и
$12,1,1$. Объясните различие между числом 12 в массиве и его
двумя символами VGA. Почему после смены максимума надо менять
проверку сохранённого значения, но не делитель в \texttt{vga\_two}?

\paragraph{45 минут: диск и повреждение копии.}
Запишите смещения \texttt{GR16}, оценок и checksum относительно
начала рассчитанного сектора. После сохранения скопируйте
\texttt{grades-lab.img} в новый \texttt{grades-corrupt.img}.
Для уже сохранённого первого байта оценки 12 измените \emph{в копии}
ровно этот байт на 1, не пересчитывая сумму:
\begin{lstlisting}[basicstyle=\ttfamily\footnotesize]
cp -n grades-lab.img grades-corrupt.img
printf '\001' | dd of=grades-corrupt.img bs=1 \
  seek="$((OFFSET + 8))" conv=notrunc
\end{lstlisting}
Предварительно положите в \texttt{OFFSET} найденное байтовое
смещение начала сектора. \texttt{printf '\textbackslash001'} выдаёт
один байт, \texttt{bs=1} задаёт единицу смещения,
\texttt{seek=} выбирает первый байт оценок, \texttt{conv=notrunc}
оставляет остальное неизменным. Запустите журнал с повреждённой
копией вместо целой, наблюдайте \texttt{BAD\_CHECKSUM}.
Измените оценку в RAM и повторите загрузку клавишей E:
плохой сектор не должен заменить текущие данные. При повторном опыте
возьмите \emph{новое} имя копии.

\paragraph{35 минут: отчёт с собственными проверками.}
Представьте изменённый \texttt{journal.asm}, результат
\texttt{diff -u} с эталоном этапа 6 и для каждого изменения
следующую запись: место в коде, собственный расчёт (основание и
единицы), старая/новая константы, выполненные команды, ожидаемое,
наблюдаемое по экрану/COM1/образу и объяснение результата.
Отдельно покажите, что сохранение оценок переживает перезапуск,
а образ повреждался только в копии.
\endgroup
